\documentclass[12pt,a4paper]{report}
\usepackage[utf8]{inputenc}
\usepackage[T1]{fontenc}
\usepackage{lmodern}
\usepackage{graphicx}
\usepackage{amsmath,amssymb}
\usepackage{hyperref}
\usepackage{geometry}
\usepackage{caption}
\usepackage{booktabs}
\usepackage{enumitem}
\usepackage{float}
\usepackage{listings}
\usepackage{color}

\geometry{
  left=30mm,
  right=25mm,
  top=25mm,
  bottom=30mm
}

\hypersetup{
  colorlinks=true,
  linkcolor=black,
  urlcolor=blue,
  citecolor=black,
  pdftitle={Apriori Algorithm and Its Implementation in a Web Application},
  pdfauthor={Student Name},
  pdfsubject={University Project}
}

\lstset{
  basicstyle=\ttfamily\small,
  breaklines=true,
  frame=single,
  numbers=left,
  numberstyle=\tiny,
  showstringspaces=false,
  keywordstyle=\color{blue},
  commentstyle=\color{gray}
}

\begin{document}

% -------------------------
% Title page
% -------------------------
\begin{titlepage}
  \centering
  {\scshape\LARGE University Name \par}
  \vspace{2cm}
  {\scshape\Large Department of Computer Science \par}
  \vspace{2.5cm}
  {\huge\bfseries Apriori Algorithm and Its Implementation in a Web Application \par}
  \vspace{1.5cm}
  {\Large\itshape Student Name\par}
  \vfill
  A project report submitted in partial fulfillment of the requirements for the\\
  degree of Bachelor/Master of Science in Computer Science

  \vspace{1.5cm}
  {\large Supervisor: Dr. Supervisor Name \par}
  \vspace{2cm}
  {\large May 1, 2026 \par}
\end{titlepage}

% -------------------------
% Abstract
% -------------------------
\begin{abstract}
This report presents the Apriori algorithm for association rule learning and describes a practical implementation integrated into a web application. The first part introduces the theoretical foundations of association rules and the key measures—support, confidence and lift. The second part details the system architecture, technology choices (Python backend and a React frontend), preprocessing steps for transforming raw datasets into transaction sets, and the Apriori implementation. The third part demonstrates results, including frequent itemset generation (L1, L2, L3), candidate generation and pruning, and the extraction and interpretation of association rules. Placeholders for screenshots and examples are included to illustrate the user interface and algorithm stages. The report concludes with lessons learned and directions for future work.
\end{abstract}

% -------------------------
% Table of contents
% -------------------------
\tableofcontents
\clearpage

% -------------------------
% General introduction
% -------------------------
\chapter*{General Introduction}
\addcontentsline{toc}{chapter}{General Introduction}

In many domains, particularly retail and e-commerce, understanding relationships between items purchased together is crucial for recommendations, promotions, and inventory management. Association rule learning is a widely-used technique for discovering interesting relations between variables in large transactional datasets. The Apriori algorithm is one of the earliest and most influential algorithms for frequent itemset mining and association rule generation.

This report aims to (1) explain the theoretical foundations of Apriori, (2) describe a web-based implementation using modern web technologies, and (3) present example outputs and visualizations that aid interpretation. The target audience is undergraduate/graduate students and developers who want a clear, reproducible guide to integrating Apriori into an application.

% -------------------------
% Chapter 1: Theoretical background
% -------------------------
\chapter{Theoretical Background}

\section{Association Rule Learning}
Association rules are implications of the form
\[
X \Rightarrow Y
\]
where \(X\) and \(Y\) are disjoint itemsets (i.e., \(X \cap Y = \varnothing\)). Rules indicate that when items in \(X\) occur in a transaction, items in \(Y\) are likely to occur as well.

Association rule learning typically consists of two steps:
\begin{enumerate}[nosep]
  \item Find all frequent itemsets whose support exceeds a user-specified minimum support threshold.
  \item Generate high-confidence rules from those frequent itemsets whose confidence is above a user-specified minimum confidence threshold.
\end{enumerate}

\subsection{Key Measures}
Three commonly used metrics are support, confidence, and lift.

Support:
\[
\text{support}(A) = \frac{|\{t \in \mathcal{T} : A \subseteq t\}|}{|\mathcal{T}|}
\]
where \(\mathcal{T}\) is the set of all transactions and \(|\cdot|\) denotes cardinality. For an itemset \(A\), support is the proportion of transactions that contain \(A\).

Confidence (for rule \(X \Rightarrow Y\)):
\[
\text{confidence}(X \Rightarrow Y) = \frac{\text{support}(X \cup Y)}{\text{support}(X)}
\]
It represents the conditional probability that transactions containing \(X\) also contain \(Y\).

Lift:
\[
\text{lift}(X \Rightarrow Y) = \frac{\text{confidence}(X \Rightarrow Y)}{\text{support}(Y)} = \frac{\text{support}(X \cup Y)}{\text{support}(X)\cdot \text{support}(Y)}
\]
Lift measures how much more often \(X\) and \(Y\) occur together than expected if they were statistically independent. A lift greater than 1 indicates positive correlation.

\section{Apriori Algorithm}
Apriori is an iterative algorithm that uses a level-wise search to identify frequent itemsets. It relies on the Apriori property: any subset of a frequent itemset must also be frequent. This property enables candidate pruning.

\subsection{High-level steps}
\begin{enumerate}[nosep]
  \item Generate \(L_1\), the set of frequent 1-itemsets (those with support \(\geq\) min\_support).
  \item For \(k = 2, 3, \ldots:\)
    \begin{enumerate}[nosep]
      \item Generate candidate set \(C_k\) by joining pairs in \(L_{k-1}\).
      \item Prune candidates in \(C_k\) that have any \((k-1)\)-subset not in \(L_{k-1}\).
      \item Count support for remaining candidates using the dataset.
      \item Form \(L_k\) by selecting candidates in \(C_k\) with support \(\geq\) min\_support.
      \item Stop when \(L_k\) is empty.
    \end{enumerate}
  \item Generate association rules from frequent itemsets: for a frequent itemset \(I\), for every non-empty subset \(A \subset I\), consider rule \(A \Rightarrow I\setminus A\) and compute its confidence; keep rules with confidence \(\geq\) min\_confidence.
\end{enumerate}

\subsection{Pseudocode}
\begin{lstlisting}[language=Python, caption={Apriori pseudocode (outline)}]
# Pseudocode outline for Apriori
def apriori(transactions, min_support):
    L1 = find_frequent_1_itemsets(transactions, min_support)
    L = [L1]
    k = 2
    while L[k-2] is not empty:
        Ck = apriori_gen(L[k-2])  # join and prune
        count_support(Ck, transactions)
        Lk = {c for c in Ck if c.support >= min_support}
        L.append(Lk)
        k += 1
    return union of all L

def generate_rules(frequent_itemsets, min_conf):
    rules = []
    for itemset in frequent_itemsets with size >= 2:
        for each non-empty subset A of itemset:
            B = itemset - A
            conf = support(itemset) / support(A)
            if conf >= min_conf:
                rules.append((A, B, conf, lift(A,B)))
    return rules
\end{lstlisting}

% -------------------------
% Chapter 2: System design and implementation
% -------------------------
\chapter{System Design and Implementation}

\section{Technologies Used}
The system uses a two-tier architecture with a backend API and a frontend web interface:

\begin{itemize}[nosep]
  \item \textbf{Backend:} Python (3.9+). Implementation can use either \textbf{FastAPI} for an asynchronous, modern API or \textbf{Flask} for a minimal synchronous server. The backend handles file uploads, preprocessing, Apriori computation, and returns results as JSON.
  \item \textbf{Frontend:} React (Vite) — a single-page application providing the dataset upload UI, previews, parameter controls (min support, min confidence), and visual displays (frequent itemsets, association rules, relationship maps).
  \item \textbf{Libraries:} pandas for data handling, numpy for numeric ops, and optionally visualization libraries such as D3.js or Chart.js in the frontend.
\end{itemize}

\section{System Architecture}
The architecture consists of two main components:

\begin{figure}[H]
  \centering
  \includegraphics[width=0.85\textwidth]{images/architecture_diagram.png}
  \caption{High-level system architecture: React frontend communicates with Python backend via REST API endpoints.}
  \label{fig:architecture}
\end{figure}

\noindent The backend exposes endpoints:
\begin{itemize}[nosep]
  \item \texttt{POST /upload} — accepts CSV dataset.
  \item \texttt{POST /apriori} — receives algorithm parameters (min\_support, min\_confidence, columns mapping), runs Apriori, returns frequent itemsets and rules.
  \item \texttt{GET /status} — optional progress/status endpoint.
\end{itemize}

\section{Data Preprocessing and Transformation}
Real-world datasets often contain categorical attributes, multiple records per transaction, or event logs. Preprocessing steps include:

\begin{enumerate}[nosep]
  \item \textbf{Column mapping:} Identify transaction ID column and item column (or multiple item columns).
  \item \textbf{Filtering:} Remove empty/NaN entries and optionally filter by date ranges or categories.
  \item \textbf{Grouping into transactions:} Group dataset by transaction ID to produce a list of transactions, each transaction being a set or list of items.
\end{enumerate}

Example transformation in pandas:
\begin{lstlisting}[language=Python, caption={Transaction transformation example using pandas}]
import pandas as pd

df = pd.read_csv('transactions.csv')
# assume columns: transaction_id, item_description
transactions = df.groupby('transaction_id')['item_description'].apply(lambda s: list(set(s))).tolist()
\end{lstlisting}

\section{Implementation of Apriori}
The backend implements Apriori in Python. Key implementation notes:

\begin{itemize}[nosep]
  \item Represent itemsets as frozenset for dictionary keys and set operations.
  \item Use a hashmap/dictionary to count supports for candidate itemsets in one pass over transactions.
  \item Candidate generation (join step) constructs \(C_k\) by combining \(L_{k-1}\) itemsets that share \(k-2\) items.
  \item Prune candidates that have any \((k-1)\)-subset not in \(L_{k-1}\).
  \item Use efficient subset testing: for each transaction \(T\), generate all \(k\)-subsets of \(T\) and increment counts for candidates appearing in that transaction (possible optimization: only check candidates directly).
\end{itemize}

A simplified Python function (conceptual):
\begin{lstlisting}[language=Python,caption={Simplified Apriori implementation (concept)}]
from itertools import combinations
from collections import defaultdict

def apriori(transactions, min_support):
    N = len(transactions)
    # Count single items
    counts = defaultdict(int)
    for t in transactions:
        for item in t:
            counts[item] += 1
    L1 = {frozenset([i]) for i,c in counts.items() if c/N >= min_support}
    L = [L1]
    k = 2
    while True:
        prev = L[-1]
        # generate candidates Ck
        candidates = set()
        prev_list = list(prev)
        for i in range(len(prev_list)):
            for j in range(i+1, len(prev_list)):
                a = prev_list[i]
                b = prev_list[j]
                union = a.union(b)
                if len(union) == k:
                    # pruning: all (k-1)-subsets must be frequent
                    valid = all(frozenset(s) in prev for s in combinations(union, k-1))
                    if valid:
                        candidates.add(union)
        # count support
        counts_k = defaultdict(int)
        for t in transactions:
            tset = set(t)
            for c in candidates:
                if c.issubset(tset):
                    counts_k[c] += 1
        Lk = {c for c,count in counts_k.items() if count/N >= min_support}
        if not Lk:
            break
        L.append(Lk)
        k += 1
    return L
\end{lstlisting}

The backend returns a JSON object with:
\begin{itemize}[nosep]
  \item \texttt{frequent\_itemsets}: list of \{itemset, support\}
  \item \texttt{association\_rules}: list of \{antecedents, consequents, support, confidence, lift\}
\end{itemize}

% -------------------------
% Chapter 3: Results and visualization
% -------------------------
\chapter{Results and Visualization}

This chapter illustrates typical outputs from the Apriori implementation. Screenshots are included as placeholders indicating what to capture -- replace the placeholder images with actual screenshots from your web app.

\section{Dataset Upload and Preview}
\begin{figure}[H]
  \centering
  \includegraphics[width=0.9\textwidth]{images/dataset_upload.png}
  \caption{Dataset upload interface (placeholder). This screenshot should show the file selection button, parameter fields (transaction ID column, item column), and the min support/min confidence controls.}
  \label{fig:upload}
\end{figure}

\begin{figure}[H]
  \centering
  \includegraphics[width=0.9\textwidth]{images/dataset_preview.png}
  \caption{Dataset preview (placeholder). Show a few rows of the raw CSV so the user can confirm column mappings.}
  \label{fig:preview}
\end{figure}

\section{Transaction Transformation}
\begin{figure}[H]
  \centering
  \includegraphics[width=0.9\textwidth]{images/transaction_transformation.png}
  \caption{Transaction transformation (placeholder). Show the resulting transactions produced from grouping by transaction id, typically a small table: transaction id \\\rightarrow item list.}
  \label{fig:transformation}
\end{figure}

\section{Frequent Itemsets: L1, C2, L2, ...}
We demonstrate the iterative generation of frequent itemsets. For a small toy dataset, the process is:

\begin{table}[H]
  \centering
  \caption{Toy dataset (example).}
  \begin{tabular}{ll}
    \toprule
    Transaction ID & Items \\
    \midrule
    T1 & \{bread, milk\} \\
    T2 & \{bread, diaper, beer, eggs\} \\
    T3 & \{milk, diaper, beer, coke\} \\
    T4 & \{bread, milk, diaper, beer\} \\
    T5 & \{bread, milk, diaper, coke\} \\
    \bottomrule
  \end{tabular}
\end{table}

Assume min\_support = 0.6 (i.e., itemset must appear in at least 3 of 5 transactions).

\subsection{L1 (frequent 1-itemsets)}
Calculate support for each item:
\begin{align*}
  \text{support}(\text{bread}) &= 4/5 = 0.8 \\
  \text{support}(\text{milk}) &= 4/5 = 0.8 \\
  \text{support}(\text{diaper}) &= 4/5 = 0.8 \\
  \text{support}(\text{beer}) &= 3/5 = 0.6 \\
  \text{support}(\text{coke}) &= 2/5 = 0.4 \\
  \text{support}(\text{eggs}) &= 1/5 = 0.2
\end{align*}

L1 = \{bread, milk, diaper, beer\}.

\begin{figure}[H]
  \centering
  \includegraphics[width=0.9\textwidth]{images/L1_table.png}
  \caption{L1 table (placeholder). Include item and support columns for frequent 1-itemsets.}
  \label{fig:L1}
\end{figure}

\subsection{Candidate generation: C2}
Join step for \(L_1\) yields candidate 2-itemsets \(C_2\): \{bread,milk\}, \{bread,diaper\}, \{bread,beer\}, \{milk,diaper\}, \{milk,beer\}, \{diaper,beer\}. Count supports across transactions, prune those below min\_support.

\begin{figure}[H]
  \centering
  \includegraphics[width=0.9\textwidth]{images/C2_counts.png}
  \caption{C2 candidate counts (placeholder). This screenshot should show candidate pairs and their support counts.}
  \label{fig:C2}
\end{figure}

\subsection{L2 (frequent 2-itemsets)}
After pruning:
\[
L_2 = \{\{\text{bread,milk}\}, \{\text{bread,diaper}\}, \{\text{milk,diaper}\}, \{\text{diaper,beer}\}\}
\]
(Example values depend on dataset.)

\begin{figure}[H]
  \centering
  \includegraphics[width=0.9\textwidth]{images/L2_table.png}
  \caption{L2 table (placeholder). Frequent 2-itemsets and their support values.}
  \label{fig:L2}
\end{figure}

\section{Pruning Examples}
Pruning eliminates candidate itemsets having a subset not present in previous frequent sets. The following screenshot should illustrate an example where a candidate is pruned because one of its 1-item subsets was not frequent.

\begin{figure}[H]
  \centering
  \includegraphics[width=0.9\textwidth]{images/pruning_example.png}
  \caption{Pruning example (placeholder). Show candidate and the missing subset used to prune it.}
  \label{fig:pruning}
\end{figure}

\section{Association Rules and Interpretation}
From frequent itemsets we generate candidate rules. Example: from itemset \{diaper, beer\} with support \(0.6\), we can propose rules:

\begin{itemize}[nosep]
  \item \(\{\text{diaper}\} \Rightarrow \{\text{beer}\}\) with
  \[
  \text{confidence} = \frac{\text{support}(\{\text{diaper,beer}\})}{\text{support}(\{\text{diaper}\})}.
  \]
  \item \(\{\text{beer}\} \Rightarrow \{\text{diaper}\}\) similarly.
\end{itemize}

Compute lift to gauge strength beyond random occurrence.

\begin{figure}[H]
  \centering
  \includegraphics[width=0.95\textwidth]{images/final_rules_table.png}
  \caption{Final association rules table (placeholder). Include antecedent, consequent, support, confidence, lift, and short interpretation/insight per rule.}
  \label{fig:final-rules}
\end{figure}

\section{Visualization}
Visualizations help to interpret rules. Examples include:
\begin{itemize}[nosep]
  \item Bar charts of top frequent itemsets by support.
  \item Network/relationship maps (nodes = items, edges = rules with weight = confidence/lift).
  \item Interactive rule cards with buttons to promote, discount, or create marketing bundles.
\end{itemize}

\begin{figure}[H]
  \centering
  \includegraphics[width=0.9\textwidth]{images/relationship_map.png}
  \caption{Relationship map visualization (placeholder). Shows items as nodes and rules as directed edges; edge thickness encodes confidence/lift.}
  \label{fig:relationship-map}
\end{figure}

% -------------------------
% Conclusion and future work
% -------------------------
\chapter*{Conclusion}
\addcontentsline{toc}{chapter}{Conclusion}

This project demonstrated the theory and practice of association rule mining using Apriori and integrated the algorithm into a web application enabling dataset upload, parameter selection, and interactive visualizations. Apriori is an effective, conceptually simple approach for small to moderately-sized datasets, though its candidate-generation approach can become expensive on large or dense datasets.

Key takeaways:
\begin{itemize}[nosep]
  \item Apriori relies on the anti-monotonicity (Apriori) property to prune search space.
  \item Practical implementations require thoughtful preprocessing and efficient counting strategies (hash maps, transactional indexing).
  \item A web interface significantly improves accessibility for analysts, enabling instant exploration of results.
\end{itemize}

\section*{Future Work}
Potential improvements and extensions:
\begin{itemize}[nosep]
  \item Implement more scalable algorithms such as FP-Growth to handle larger datasets without expensive candidate generation.
  \item Add incremental mining to update frequent itemsets when new transactions arrive.
  \item Improve UI/UX: add interactive filtering, threshold sliders, export of rules, and A/B testing workflows for recommended bundles.
  \item Add better performance optimizations in backend (Cython, vectorized counting, parallel processing).
  \item Evaluate rules using business KPIs (lift vs. realized uplift in sales) by integrating historical promotion data.
\end{itemize}

% -------------------------
% References
% -------------------------
\begin{thebibliography}{9}
\bibitem{agrawal1994}
R. Agrawal and R. Srikant.
\newblock Fast algorithms for mining association rules.
\newblock In \emph{Proceedings of the 20th International Conference on Very Large Data Bases (VLDB)}, 1994.

\bibitem{han2011}
Jiawei Han, Micheline Kamber, and Jian Pei.
\newblock \emph{Data Mining: Concepts and Techniques (3rd Edition)}.
\newblock Morgan Kaufmann, 2011.

\bibitem{fpgrowth}
Han, J. et al.
\newblock Mining frequent patterns without candidate generation (FP-Growth).
\newblock \emph{ACM SIGMOD Record}, 2000.

\bibitem{fastapi}
Sebastian Ram\'irez.
\newblock FastAPI: Modern, fast (high-performance) web framework for building APIs with Python 3.6+.
\newblock \url{https://fastapi.tiangolo.com/}

\bibitem{react}
Facebook Inc.
\newblock React — A JavaScript library for building user interfaces.
\newblock \url{https://reactjs.org/}
\end{thebibliography}

% -------------------------
% Appendix: Additional material and screenshot placeholders
% -------------------------
\appendix
\chapter{Appendix: Screenshot placeholders and notes}

Below is the list of recommended screenshots to capture and include in the final PDF. Replace the placeholder image files (in \texttt{images/}) with actual captures:

\begin{enumerate}[nosep]
  \item \textbf{Dataset upload interface:} file chooser, parameter fields, min support/confidence sliders. (Replace \texttt{images/dataset\_upload.png}.)
  \item \textbf{Dataset preview:} show a tabular view of raw CSV columns and sample rows. (Replace \texttt{images/dataset\_preview.png}.)
  \item \textbf{Transaction transformation:} grouping and resulting transactions list or preview. (Replace \texttt{images/transaction\_transformation.png}.)
  \item \textbf{L1 table:} frequent 1-itemset table. (Replace \texttt{images/L1\_table.png}.)
  \item \textbf{C2 and counting step:} candidate list with counts. (Replace \texttt{images/C2\_counts.png}.)
  \item \textbf{L2 table and pruning example:} show pruning logic removing candidates. (Replace \texttt{images/L2\_table.png} and \texttt{images/pruning\_example.png}.)
  \item \textbf{Final rules table:} antecedent, consequent, support, confidence, lift, interpretation. (Replace \texttt{images/final\_rules\_table.png}.)
  \item \textbf{Relationship map:} network visualization of items and rules. (Replace \texttt{images/relationship\_map.png}.)
\end{enumerate}

\bigskip
Notes on capturing screenshots:
\begin{itemize}[nosep]
  \item Use a consistent viewport (e.g., 1280×720) for clarity.
  \item Highlight the UI controls used (e.g., active min\_support value).
  \item If possible, include a short caption describing the purpose of each screenshot.
\end{itemize}

\end{document}
